\documentclass[10pt,a4paper]{amsart}
\usepackage[margin=19mm]{geometry}
\usepackage{amsmath,amssymb,mathtools}
\usepackage[scr=rsfs,cal=euler]{mathalfa}
\usepackage{xfrac}
\usepackage[unicode,hidelinks,bookmarks=false]{hyperref}
\newcommand{\ie}{i.e.\ }
\newcommand{\cf}{cf.\ }
\newcommand{\resp}{respectively\ }
\newcommand{\Subsection}{\subsection}
\providecommand{\nobreakdash}{\nobreak\hbox{-}}
\setlength{\emergencystretch}{2em}
\multlinegap0pt
\usepackage{latexsym,mathrsfs}
\usepackage{units}
\usepackage{esint}
\numberwithin{equation}{section}
\usepackage{url}
\setlength{\oddsidemargin}{1pt}
\setlength{\evensidemargin}{1pt}
\setlength{\topmargin}{1pt}       
\setlength{\textheight}{650pt}    
\setlength{\textwidth}{460pt}     
\def\dsp{\def\baselinestretch{1.37}\large
\normalsize}
\newtheorem{definition}{Definition}
\newtheorem{claim}{Claim}
\newtheorem{exercise}{Exercise}
\newtheorem{example}{Example}
\newtheorem{lemma}{Lemma}[section]
\newtheorem{prop}{Proposition}
\renewcommand{\d}{\,{\rm d}}
\def\MM{\mathfrak{M}}
\def\BB{\mathfrak{B}}
\def\CC{\mathfrak{C}}
\def\leb{{\mathcal L}}
\DeclareMathOperator{\m}{\mathfrak{M}}
\DeclareMathOperator{\B}{\mathfrak{B}}
\DeclareMathOperator{\sgn}{\operatorname{sgn}}
\renewcommand{\Re}{\text{Re}}
\renewcommand{\Im}{\text{Im}}
\renewcommand{\L}{\mathcal{L}}
\DeclareMathOperator{\C}{\mathbb{C}}
\DeclareMathOperator{\Sb}{\mathbb{S}}
\DeclareMathOperator{\R}{\mathbb{R}}
\DeclareMathOperator{\Q}{\mathbb{Q}}
\DeclareMathOperator{\N}{\mathbb{N}}
\DeclareMathOperator{\Z}{\mathbb{Z}}
\DeclareMathOperator{\D}{\mathbb{D}}
\DeclareMathOperator{\T}{\mathcal{T}}
\newtheorem{theorem}{Theorem}[section]
\newtheorem{prob}{Problem}
\begin{document}
\title[An Approximate Version of Vu's Theorem on Economical Subbase]{An Approximate Version of Vu's Theorem on Economical Subbases For Non-Integer Exponents}
\author{Ataleshvara Bhargava}
\date{}
\maketitle
\noindent\emph{Source and licence.} An Approximate Version of Vu's Theorem on Economical Subbases For Non-Integer Exponents. Version arXiv:2609.27823v1 (CC BY 4.0). This material is distributed under the Creative Commons Attribution 4.0 International licence, http://creativecommons.org/licenses/by/4.0/, as recorded in the arXiv OAI licence metadata for this exact version and shown on the abstract page https://arxiv.org/abs/2609.27823v1. Full rights notice: licenses/arxiv-2609.27823-CC-BY-4.0.txt.\par\smallskip
\noindent\textit{Editorial scope.} Counted unit: the original \texttt{lemma} environment \texttt{lem:kernel} at source lines 355--366 together with its complete proof at lines 368--369; the delivered document keeps source lines 336--370 so that every referenced label and every citation resolves inside it. Native editable source is the paper's own concatenated TeX; the AMS wrapper is editorial formatting only. No publisher class, upstream script or unrelated artwork is redistributed.\par\smallskip\noindent\textit{Reference note.} This article ships no compiled bibliography (biblatex); the entries delivered here are composed from the author's own BibTeX (.bib) fields, with no field invented and no entry dropped.
\section{The Davenport-Heilbronn Method} \label{sec:DH}
This section is devoted to proving some necessary estimates on a weighted Diophantine counting problem. Our approach is based on \cite{WoolThin}. We will make use of Freeman's refined version of this method; the reader may consult \cite{FreemanLB}, \cite{FreemanAsymp}, \cite{Poulias1} or \cite{WooleyDHF} for more information on this. See also Chapter 20 in \cite{DavBook}. Recall for $0 < \Lambda \in \mathbb{R}$ the definitions $Y = \Lambda^{1/\theta}$ and $Y_s = \Lambda^{\frac{1}{s\theta}}$. 
Let $I_1 = \cdots = I_{s-1} = (Y_s,Y]$ and let $I_s = (c_0Y, Y]$, where $c_0 = \eta_1\eta_2^{-1}$. Let $\mathcal{T} = I_1\times \cdots \times I_s$ and recall that 
\[ U_{s,\delta}(\Lambda; \mathcal{T}) = \sum_{\substack{ (x_1,\ldots,x_s) \in \mathcal{T} \\ |x_1^{\theta}+\cdots+x_s^{\theta}-\Lambda|<\delta }} (x_1\cdots x_s)^{-1+\theta/s}. \]
The main result of this section is the following. 

\begin{theorem} \label{thm:DH} 
Fix $\tau > 0$. Let $\theta > 2$ and let $s \geq s_0+\chi$. There exist positive constants $ \Upsilon_1(s,\theta) $ and $\Upsilon_2(s,\theta)$ such that for all sufficiently large real $\Lambda > 0$ one has 
\[ \tau\Upsilon_1(s,\theta) \leq U_{\tau/2,s}(N,\mathcal{T})\leq \tau\Upsilon_2(s,\theta). \]
\end{theorem}
To initiate our proof of Theorem \ref{thm:DH}, define the weighted exponential sums 
\[ g_1(\alpha) = \sum_{x \in I_1} x^{-1+\theta/s}e(\alpha x^{\theta}) \quad \text{ and } \quad g_2(\alpha) = \sum_{x \in I_s} x^{-1+\theta/s}e(\alpha x^{\theta}). \]
We need a certain kernel function with some predefined properties. For $\delta > 0$, let $\tilde{\delta} = \delta\log(Y)^{-1}$. Define the kernel functions $K^{\delta}_+(\alpha)$ and $K^{\delta}_-(\alpha)$ by 
\begin{align} \label{eq:kerform}
K^{\delta}_{\pm}(\alpha) = \frac{\sin(\pi\alpha\tilde{\delta})\sin(\pi\alpha(2\delta\pm\tilde{\delta}))}{\pi^2 \alpha^2 \tilde{\delta}} = (2\delta\pm\tilde{\delta})\frac{\sin(\pi \alpha \tilde{\delta})}{\pi\alpha\tilde{\delta}} \cdot \frac{\sin(\pi\alpha(2\delta\pm \tilde{\delta}))}{\pi \alpha (2\delta\pm\tilde{\delta})}.
\end{align} 
Define the functions $\psi_+(\xi)$ and $\psi_-(\xi)$ by 
\[ \psi_{\pm}(\xi) = \int_{\mathbb{R}} e(\xi\alpha)K^{\tau}_{\pm}(\alpha)\d\alpha. \]
For measurable $B \subseteq \mathbb{R}$, let $\chi_{B}$ denote the indicator function of $B$. One has the following lemma. 

\begin{lemma} \label{lem:kernel}
Fix $\delta > 0$. The following hold: 
\begin{align} \label{eq:kerbound}
K^{\delta}_{\pm}(\alpha) \ll_{\delta} \min\{1, |\alpha|^{-1},\log(Y)|\alpha|^{-2} \},
\end{align}
\begin{align} \label{eq:kerasymp}
K^{\delta}_{\pm}(\alpha) = 2\delta+O(\log(Y)^{-1}) 
\end{align}
with \eqref{eq:kerasymp} holding for $|\alpha| < Y^{-\beta}$ for any fixed $\beta > 0$. Moreover, one has
\[ \chi_{(-\delta+\tilde{\delta},\delta-\tilde{\delta})} (\xi) \leq \psi_-(\xi) \leq \chi_{(-\delta,\delta)}(\xi) \leq \psi_+(\xi) \leq \chi_{(-\delta-\tilde{\delta},\delta+\tilde{\delta})}(\xi). \]
Further, $|\psi_{\pm}(\xi)-\chi_{(-\delta,\delta)}(\xi)| = 0$ whenever $||\xi|-\delta|>\tilde{\delta}$ and is at most $1$ when $||\xi|-\delta|\leq \tilde{\delta}$. 
\end{lemma}

\begin{proof} See Lemma 3.1 and the discussion following it in \cite{Bharloc}. See also Lemma 1 in \cite{FreemanAsymp}, Lemma 2.1 in \cite{Poulias1} and Section 3.2.1 in \cite{Poulthesis}. 
\end{proof}

\begin{thebibliography}{99}
\bibitem[]{Bharloc} Bhargava, A.. Sharp Transitions for Localized Solutions to a {D}iophantine Inequality. (2026+)
\bibitem[]{DavBook} Davenport, H.. Analytic methods for {D}iophantine equations and {D}iophantine inequalities. (2005) xx+140
\bibitem[]{FreemanAsymp} Freeman, D. E.. Asymptotic lower bounds and formulas for {D}iophantine inequalities. Number theory for the millennium, {II} ({U}rbana, {IL}, 2000) (2002) 57--74
\bibitem[]{FreemanLB} Freeman, D. E.. Asymptotic lower bounds for {D}iophantine inequalities. Mathematika 47 (2000) 127--159
\bibitem[]{Poulias1} Poulias, C.. Diophantine inequalities of fractional degree. Mathematika 67 (2021) 949--980
\bibitem[]{Poulthesis} Poulias, K.. On {D}iophantine problems involving fractional powers of integers. (2021)
\bibitem[]{WoolThin} Wooley, T. D.. On {V}u's thin basis theorem in {W}aring's problem. Duke Math. J. 120 (2003) 1--34
\bibitem[]{WooleyDHF} Wooley, T. D.. On {D}iophantine inequalities: {F}reeman's asymptotic formulae. Proceedings of the {S}ession in {A}nalytic {N}umber {T}heory and {D}iophantine {E}quations 360 (2003) 1--32
\end{thebibliography}
\end{document}
